

Paragraph 3 theme: Distinguishing blink artifacts from blink detection in the literature.
It is important to distinguish between two major perspectives on blink activity in the existing literature. In many EEG studies, blink-related activity is treated as an artifact because it contaminates neural signals and may interfere with the interpretation of brain activity. In such cases, blink components are commonly removed during EEG preprocessing. However, in fatigue-driving studies, blink activity can also be treated as meaningful information rather than noise. From this perspective, blink events are identified, extracted, and analysed as potential indicators of driver fatigue. Blink detection methods can generally be classified into machine-learning-based, morphology-based, and thresholding-based approaches. Machine learning methods rely on trained models to classify blink events, while morphology-based methods identify blinks based on their characteristic waveform shape. Thresholding-based methods, in contrast, detect blinks by comparing signal amplitudes with predefined or adaptive threshold values.